\documentclass[11pt]{article}
\usepackage{setspace}
\usepackage{geometry}
\geometry{margin=1in}
\setstretch{1.15}

\begin{document}

\noindent
\textbf{Zhongli Wang} \\
Phone: +49 160 9541 6015 \\
Email: nkwangzhongli@gmail.com \\
Website: https://sites.google.com/view/zhongli-wang \\
Bielefeld, Germany \\

\vspace{0.5cm}

\noindent
\textbf{Professor Topi Miettinen and the Recruitment Committee} \\
Department of Finance and Economics \\
Hanken School of Economics \\
Helsinki, Finland \\

\vspace{0.5cm}

\noindent
\textbf{Application for the Postdoctoral Researcher Position in Economics} \\

\vspace{0.3cm}

\noindent Dear Professor Miettinen and Members of the Recruitment Committee,\\[0.2cm]

I am writing to apply for the postdoctoral researcher position in economics at Hanken School of Economics. I am currently completing my Ph.D.\ in Economics at Bielefeld University and Università Cattolica del Sacro Cuore within the Marie Skłodowska--Curie EPOC programme, and I expect to finish before the start of the position. My research combines industrial organization, behavioral economics, and computational methods to study firm behavior, market dynamics, and decision-making under frictions and heterogeneity. Hanken's strong research orientation, its close integration with Helsinki GSE, and its particular strengths in Industrial Organization, Behavioral Economics, and Microeconomics make this position especially attractive to me.

My doctoral work develops computational and empirical models of strategic interaction in technology-intensive markets. In my dissertation, \textit{Strategic Behavior and Economic Dynamics in ICT Markets: An Agent-based Framework}, I study how firms and heterogeneous consumers interact under network effects, congestion, and information frictions. The broader goal of this work is to understand how incentives, market design, and behavioral responses shape market outcomes and welfare. Although my main empirical setting is digital markets, the underlying themes of my research (strategic adaptation, bounded rationality, and policy-relevant counterfactual analysis) connect closely to Hanken's focus areas in industrial organization, behavioral economics, and the economic analysis of innovation and entrepreneurship.

Methodologically, I combine structural reasoning with modern quantitative tools. My work relies on agent-based modeling, simulation-based econometrics, and machine-learning-assisted approximation methods to analyze settings in which analytical solutions alone are insufficient. At the same time, I aim to keep the research grounded in empirically meaningful questions and transparent identification strategies. This combination of theory, computation, and data is also reflected in my collaborative projects. In \textit{Polytope Fraud Theory} (\textit{International Review of Financial Analysis}, 2025), I worked on supervised learning approaches for detecting structured irregularities in firm financial statements. In \textit{Agentic Workflows for Economic Research}, I helped develop reproducible AI-assisted research pipelines. Together, these projects demonstrate my interest in rigorous, data-driven economic research and my ability to work across methodological boundaries.

I would be excited to contribute to Hanken's research and teaching environment in several ways. First, I believe my research profile would fit particularly well with the Industrial Organization, Behavioral Economics, and Microeconomics groups within Helsinki GSE. Second, I would welcome the opportunity to engage with Hanken's behavioral and experimental infrastructure, including Helsinki LABBET, and to further develop research on economic decision-making in controlled environments. Third, I would be glad to contribute to teaching in central areas of economics. My teaching experience includes two years as a helping assistant at the University of Zurich, where I supported quantitative courses and contributed to the SYLVA automated scoring and feedback system. I enjoy teaching and see it as a natural complement to research, especially in an intellectually ambitious environment such as Hanken.

What especially attracts me to this position is the opportunity to develop an independent research profile while being embedded in a vibrant economics community with strong links to business and society. I value collaborative research environments, seminar culture, and interdisciplinary exchange, and I would be very happy to contribute actively to the department and the wider Helsinki GSE community.

Thank you so much for considering my application! I would be delighted to discuss my candidacy further and to contribute to Hanken's research and teaching in economics.

\vspace{6mm}

\noindent Sincerely,\\[12mm]
\noindent Zhongli Wang

\end{document}